\section{从互联网浪潮到 AI 浪潮：基础设施只是必要条件}

讲者把今天的 AI 与早期互联网作类比：初期资本和注意力集中在路由器、交换机、accelerator 与数据中心；真正改变生产率的阶段，则发生在 telemedicine、education、GovTech 和企业流程开始重构之后。本节因此不问“算力是否重要”，而问“算力怎样穿过组织边界，变成稳定服务和可测结果”。

\subsection{四步扩散链：capacity、application、organization、outcome}

本节把上一段历史类比转换成可验收的系统链。基础设施扩散可以分为四步：第一步提供技术容量；第二步出现具体行业 application trigger；第三步组织修改数据流、权限和工作方式；第四步才产生服务质量、生产率或福利结果。读图时先找每一步的 owner、输入和指标；任一步断裂，硬件投资都可能停留在低利用率、演示项目或不可持续补贴。

\lecturefigure{01-infrastructure-diffusion.png}{基础设施只有经过应用触发和组织重构，才形成可测量结果。}{本地字幕 03:40--06:20；概念重绘。}

\begin{knowledgebox}{读图：先找断点，不要直接连 GDP}
读图时从右向左问：目标 outcome 用什么指标衡量？哪个 workflow 真正改变？应用依赖什么数据、人才和采购？最后才问需要多少 accelerator。若一份方案只描述左侧 capacity，却没有右侧的 adoption、feedback 和 retirement rule，它仍是资源计划，不是转型计划。
\end{knowledgebox}

\begin{importantbox}{基础设施扩散的最小验收式}
可把项目价值粗略拆为
\[
V \approx C\times U\times A\times Q,
\]
其中 $C$ 是可用 capacity，$U$ 是有效利用率，$A$ 是 adoption，$Q$ 是单位使用带来的质量或生产率改进。任何一项接近零，单纯扩大 $C$ 都难以创造结果。该式不用于精确估值，而用于迫使方案补齐中间环节。
\end{importantbox}

\begin{warningbox}{历史类比不能替代当代约束}
互联网和 AI 都经历 infrastructure-first 阶段，但 AI 还受到训练数据、模型更新、推理成本、内容安全、自动化权限和 human accountability 的约束。类比的价值在于提醒“应用扩散晚于硬件热潮”，不能据此推出两条产业曲线会以相同速度或相同赢家重演。
\end{warningbox}

\subsection{本章小结}

国家级 AI 不能只用购置量、MW 或模型参数衡量。可迁移的方法是建立从 capacity 到 outcome 的完整链路，并为 utilization、adoption、quality 与退出机制分别设置指标。

